\section{Related Work}
\label{sec:related}

\paragraph{Data provenance}
Database research distinguishes \emph{why}-provenance (which input tuples justify an output),
\emph{where}-provenance (which input locations a value was copied from) and
\emph{how}-provenance (how inputs were combined)~\citep{buneman2001why,cheney2009provenance};
provenance semirings unify how- and why-provenance as annotations computed alongside a
query~\citep{green2007semirings}. Lineage tracing for data warehouses goes back to Cui, Widom and
Wiener~\citep{cui2000tracing,cui2003lineage}, and systems such as Perm compute provenance by
rewriting queries~\citep{glavic2009perm}. These are mostly \emph{dynamic}: they annotate data as
queries run. Müller and Grust derive where- and why-provenance for SQL by a value-less abstract
interpretation after recording the query's control flow~\citep{muller2015abstract}, which is
closest in spirit to Konstruct's static analysis. Cheney, Ahmed and Acar cast provenance as
dependency analysis: an output depends on an input if changing the input can change the
output~\citep{cheney2007dependency}. That definition is the one \cmd{konstruct verify} tests.
Perturbation is an established tool for explaining query results: Walenz and Yang analyse how a
query's answers respond to perturbations of its parameters~\citep{walenz2016perturbation}, and
counterfactual causality and influence-based explanations remove or change tuples to see what an
answer depends on~\citep{meliou2010causality,wu2013scorpion}. Konstruct applies the idea to a
different question: not why a particular answer is what it is, but whether a static lineage claim
about a query holds on the data.

\paragraph{Column lineage tools}
OpenLineage is an open standard for lineage events emitted by running jobs, with column-level
lineage as a dataset facet, and Marquez is its reference
implementation~\citep{openlineage,marquez}. SQLLineage and sqlglot's lineage module trace columns
through a SQL statement and run locally; sqlglot's lineage nodes carry the expression of each
step~\citep{sqllineage,sqlglot}, and Konstruct builds on sqlglot. SQLMesh, also built on sqlglot,
derives column-level lineage locally from its own understanding of a project's
SQL~\citep{sqlmesh}. dbt Explorer offers column-level lineage for dbt projects in dbt Labs' hosted
service~\citep{dbtcll}, and metadata platforms such as DataHub and OpenMetadata collect lineage
from many systems into a central service~\citep{datahub,openmetadata}. Konstruct's emphasis is
different in three ways. It joins lineage across dbt, orchestration code, dashboards and loaders
in one graph built from a repository with no service. It treats lineage as a claim to be tested
against the data rather than a fact to be displayed. And it composes each column's expression
through every model upstream, so lineage carries formulas and conditions, not only edges. It
writes OpenLineage events, so it complements rather than replaces these systems.

\paragraph{Semantic layers and Apache Ossie}
Semantic layers such as the dbt Semantic Layer with MetricFlow and Looker's LookML define metrics
once over physical tables~\citep{dbtsemantic,lookml}. Apache Ossie, formerly the Open Semantic
Interchange, is a vendor-neutral YAML specification for such models (datasets, fields,
relationships and metrics), with a portability ladder that grades how well a converter preserves
each element~\citep{ossie}. The specification describes what a metric is, but not where its
inputs come from, who owns it or whether its lineage holds. Konstruct reads Ossie models as an
input and writes its ontology as one, adding exactly those facts in a custom extension
(\cref{sec:ossie}).

\paragraph{Knowledge graphs and text-to-SQL}
Benchmarks of language models on databases show a large gap between toy and enterprise settings.
Spider measured cross-domain text-to-SQL~\citep{yu2018spider}; BIRD added large, dirty databases
and external knowledge~\citep{li2023bird}; Spider 2.0 moved to enterprise workflows across
dialects and dbt projects~\citep{lei2025spider2}. Sequeda et al.\ showed that a knowledge graph of
the business more than triples GPT-4's accuracy on enterprise questions~\citep{sequeda2023benchmark}.
Konstruct produces such a graph from the code that already exists, with entities, measures,
relationships and verified formulas, and serves it to agents through the Model Context
Protocol~\citep{mcp}.

\paragraph{Static analysis of orchestration code}
Airflow and Dagster discover pipelines by importing user Python modules~\citep{airflow,dagster}.
Importing executes arbitrary code at analysis time, which is unacceptable in a review tool that
runs on every pull request. Konstruct parses DAG and asset files with Python's \cmd{ast} module and
never imports them; its parsers are tested for parity against Airflow's own DagBag and Dagster's
own asset and job graphs on the fixtures in its test suite.

\paragraph{Testing analyses with data}
Differential testing compares two implementations on generated inputs~\citep{mckeeman1998differential},
and has found many bugs in database engines~\citep{rigger2020pqs}; property-based testing generates
the inputs~\citep{maciver2019hypothesis}. Mutation testing measures a test's power by seeding
faults and counting those it detects~\citep{jia2011mutation}. Konstruct uses all three to test its
own verifier (\cref{sec:evaluation}).
